% =============================================================================
% 6.2 Answer to the Main Research Question
% =============================================================================

\section{Answer to the Main Research Question}
\label{sec:answer}

\begin{quote}
  \textit{How do reliability, system stability, and execution performance change when a workflow orchestration system is migrated from a single-VM DockerRunLauncher to a Kubernetes Run Launcher under increasing concurrent workload, and at what specific workload level does this migration become net beneficial?}
\end{quote}

\paragraph{Direct answer.}

\textbf{Reliability} changes dramatically: the VM's success rate collapses from 100\,\% at L1--L3 to 66.7\,\% at L4, 61.9\,\% at L5, and 0.0\,\% at L6, driven by memory exhaustion (Linux OOM kills) on the shared 4\,GB host.
Kubernetes, by contrast, maintained 100\,\% success rate across all tested concurrency levels through hard per-pod memory limits (2\,GiB per run pod, accommodating the 400\,MB workload buffer and $\approx$500--600\,MB Python/Dagster overhead) that prevent cross-job memory interference.

\textbf{System stability}, measured as execution time variance, improves on Kubernetes: consistent per-pod isolation produces near-constant execution time even at L6, whereas VM variance grows as memory contention and OOM events introduce randomness into survivor timing.

\textbf{Execution performance} at L1--L3 is \emph{slower on K8s}: K8s adds 19--22\,s per-run startup overhead (pod creation + Python initialisation + PostgreSQL connection), equivalent to $\approx$30\% of the 63-second job execution time.
At higher concurrency, this overhead grows to 35--57\,s due to concurrent initialisation contention, but the VM's reliability collapse from L4 onward makes it definitively inferior.
For \emph{surviving} VM jobs at L4+, execution time explodes (168.6\,s at L4, 464.1\,s at L5) due to extreme memory contention among survivors.

\textbf{The migration becomes net beneficial at 5 concurrent Dagster jobs} on hardware comparable to the test configuration (4\,vCPU, 4\,GB RAM VM / single-node Kind cluster).
This is the \emph{reliability} crossover: the first concurrency level at which the VM cannot guarantee $\geq$95\,\% success while Kubernetes can.
Even if raw execution time for surviving VM jobs were faster at low concurrency, scheduling workloads that produce failures at 66.7\,\% success rate is operationally unacceptable for production pipelines.

Teams whose main concern is execution time stability---not just success rate---benefit from Kubernetes at any concurrency level above $\approx$3 concurrent jobs, since K8s execution time variance is lower at every tested level.

